\documentclass[12pt]{article}
\usepackage[T2A]{fontenc}
\usepackage[utf8]{inputenc}
\usepackage[russian]{babel}
\usepackage{amsmath, amssymb, amsfonts}
\usepackage{geometry}
\usepackage{fancyhdr}
\usepackage{microtype}
\usepackage{titlesec} % Пакет для управления стилем заголовков
\usepackage{enumitem}
\usepackage{tikz}
\usepackage{tikz-3dplot}
\usetikzlibrary{calc, intersections, through, backgrounds, arrows.meta}

% Настройка Section: по центру, жирный, крупный
\pagestyle{fancy}
% Настройка полей
\geometry{a5paper, margin=1cm, top=2cm, headheight=15pt}

% --- ЖЕСТКАЯ КОРРЕКЦИЯ ЦЕНТРИРОВАНИЯ ---
\pagestyle{fancy}
\fancyhf{} % Полная очистка всех полей (убирает невидимые отступы)
\fancyhead[C]{\thepage} % Номер строго по центру
\renewcommand{\headrulewidth}{0pt} % Убираем линию
\fancyfoot{} % Низ пустой

% Убираем внутренние отступы колонтитулов, которые могут давать перекос
\setlength{\headwidth}{\textwidth} 

% Переопределяем системный стиль plain, чтобы он не сбивал настройки
\makeatletter
\let\ps@plain\ps@fancy
\makeatother

\parindent=0pt
\parskip=0.6em
\titleformat{\section}
  {\normalfont\bfseries\centering}{}{0pt}{}

% Настройка Subsection: отступ 1см, жирный, обычный размер
% {команда}[формат]{стиль}{метка}{отступ_от_метки}{перед_заголовком}[после_заголовка]
\titleformat{\subsection}
  {\normalfont\bfseries}{\hspace{1cm}}{1pt}{}
\titlespacing*{\subsection}{1cm}{2ex plus 1ex minus .2ex}{0.5ex plus .2ex}

% Убираем автоматическую нумерацию, чтобы заголовки были как в оригинале
\setcounter{secnumdepth}{0}
\AddEnumerateCounter{\asbuk}{\russian@asbuk@alph}{щ} % section 8.1
\setlist{wide, nosep, font=\rmfamily\normalfont}
\setlist[enumerate, 2]{label=\asbuk*)}
\setlist[enumerate, 3]{label=\asbuk*)}


\begin{document}

\section{2.1 Введение в многомерный мир}

\begin{enumerate}
\item Сколько вершин, ребер, двумерных и трехмерных граней имеет четырехмерный куб?

\item Описать правильный четырехмерный симплекс системой неравенств.

\item Сколько вершин, ребер, двумерных и трехмерных граней имеет четырехмерный симплекс?

\item Четырехмерным правильным октаэдром называется выпуклая оболочка единичных векторов $\pm e_1 = (\pm 1, 0, 0, 0)$, \\ 
$\pm e_2 = (0, \pm 1, 0, 0)$, 
$\pm e_3 = (0, 0, \pm 1, 0)$, 
$\pm e_4 = (0, 0, 0, \pm 1)$. Сколько вершин, ребер, двумерных и трехмерных граней имеет этот многогранник?

\item Вычислить для симплекса, куба и октаэдра альтернированную сумму $v_0 - v_1 + v_2 - v_3,$ где $v_i$ -- число $i$-мерных граней.

\item Построить двумерную плоскость, пересекающую четырехмерный симплекс, 
но не пересекающую ни одно из его ребер.

\item Как вылезти «через четвертое измерение» из обычной двумерной сферы наружу 
(как соединить точку $(0,0,0,0)$ с точкой $(2,0,0,0)$ непрерывным путем, 
не пересекающим сферы $x_1^2 + x_2^2 + x_3^2 = 1$, $x_4 = 0$)?

\item В $\left(\sum\limits_{i = 1}^n a_i^2 \right)^\frac{1}{2}\mathbb{R}^4$ заданы две сцепленные окружности:
\[
O_1 = \{ x \in \mathbb{R}^4 \mid x_1^2 + x_2^2 = 1,\ x_3 = x_4 = 0 \},
\]
\[
O_2 = \{ x \in \mathbb{R}^4 \mid (x_1 - 1)^2 + x_2^2 = 1,\ x_3 = x_4 = 0 \}.
\]
Можно ли их расцепить (т.е. можно ли непрерывно перевести $O_2$, скажем, в окружность 
$(x_1 - 3)^2 + x_2^2 = 1$, не пересекая $O_1$)?

\item Правда ли, что в четырехмерном пространстве можно развязать любой узел?

\item Построить в четырехмерном пространстве две двумерные плоскости, пересекающиеся в одной точке.

\item Построить в четырехмерном пространстве две «скрещивающиеся» двумерные плоскости 
(не пересекающиеся и не параллельные).

\item Рассмотрим в четырехмерном пространстве куб с единичным ребром и сечения его 
трехмерными плоскостями, ортогональными его наибольшей диагонали. 
\begin{enumerate}
	\item Какова форма получающихся сечений?  
	\item Как зависит объем этих сечений от их расстояния до центра куба?
\end{enumerate}
\item При каких $n$ в $\mathbb{R}^4$ можно построить многогранник с $n$ вершинами, 
любая пара вершин которого связана ребром?

\item Сколько вершин, ребер, двумерных, …, $k$-мерных, …, $n-1$-мерных граней 
(в совокупности) имеет $n$-мерный куб 
\[
\Pi^n := \{ x = (x_1,\dots,x_n) \mid |x_i| \le 1,\ 1 \le i \le n \}.
\]

\item Рассмотрим $n$-мерный параллелепипед 
\[
	\Pi^n(a) := \{ x = (x_1,\dots,x_n) \mid \left| \frac{x_i}{a_i} \right| \le 1,\ 1 \le i \le n \}.
\]
Каково расстояние от начала координат до $k$-мерной грани $\Pi^n(a)$?  
Каково наибольшее количество различных чисел среди всех возможных расстояний 
от начала координат до всех граней $\Pi^n(a)$ (если варьировать $a$)?

\item Описать поляру параллелограмма $\Pi^n(a)$.

\item Найти радиус описанного шара вокруг $\Pi^n(a)$.

\item Найти радиус вписанного шара в $\Pi^n(a)$.

\item Найти радиус наибольшего круга, вписанного в $\Pi^n(a)$.

\item Найти радиус наименьшего цилиндра, описанного около $\Pi^n(a)$.

\item Сколько вершин, ребер, двумерных, …, $k$-мерных, …, $n-1$-мерных граней 
(в совокупности) имеет $n$-мерный правильный октаэдр?

\item Рассмотрим $n$-мерный октаэдр 
\[
O^n(a) := \{ x = (x_1,\dots,x_n) \mid \sum_{i=1}^n \left| \frac{x_i}{a_i} \right| \le 1 \}.
\]
Каково расстояние от начала координат до $k$-мерной грани $O^n(a)$?  
Каково наибольшее количество различных чисел среди всех возможных расстояний 
от начала координат до всех граней $O^n(a)$ (если варьировать $a$)?

\item Описать поляру октаэдра $O^n(a)$.

\item Найти радиус описанного шара вокруг $O^n(a)$.

\item Найти радиус вписанного шара в $O^n(a)$.

\item Найти радиус наибольшего круга, вписанного в $O^n(a)$.

\item Найти радиус наименьшего цилиндра, описанного около $O^n(a)$.

\item Доказать, что объём симплекса равен $\frac{1}{n} h V_{n-1}$, где $V_{n-1}$ -- $(n-1)$-мерный объём грани, $h$ -- высота, опущенная на эту грань.

\item Найдите объём $V_n$ правильного симплекса $\sum x_i = 1$, $x_i \ge 0$ ($i=1,\dots,n+1$).

\item Найти объём правильного $n$-мерного симплекса с ребром 1.

\item Пусть $V$ -- объём параллелепипеда, натянутого на векторы $v_1,\dots,v_n$. Доказать, что 
\[
V^2 = \det \bigl( \langle v_i, v_j \rangle \bigr)_{1\le i,j\le n}.
\]

\item Пусть $V$ -- объём симплекса с вершинами $A_1,\dots,A_{n+1}$, $d_{ij} = |A_i A_j|$. Доказать, что 
\[
V^2 = \left(\frac{1}{n!}\right)^2 \frac{(-1)^{n+1}}{2^n} \det \begin{pmatrix}
0 & d_{12}^2 & \dots & d_{1,n+1}^2 & 1 \\
d_{21}^2 & 0 & \dots & d_{2,n+1}^2 & 1 \\
\vdots & \vdots & \ddots & \vdots & \vdots \\
d_{n+1,1}^2 & d_{n+1,2}^2 & \dots & 0 & 1 \\
1 & 1 & \dots & 1 & 0
\end{pmatrix}.
\]

\item Пусть $R$ -- радиус описанной сферы симплекса с вершинами $A_1,\dots,A_{n+1}$, $d_{ij} = |A_i A_j|$, $V$ -- объём симплекса, 
\[
\Delta = \det \begin{pmatrix}
0 & d_{12}^2 & \dots & d_{1,n+1}^2 \\
d_{21}^2 & 0 & \dots & d_{2,n+1}^2 \\
\vdots & \vdots & \ddots & \vdots \\
d_{n+1,1}^2 & d_{n+1,2}^2 & \dots & 0
\end{pmatrix}.
\]
Доказать, что $2 R^2 V^2 + \Delta = 0$.

\item Пусть сферы радиусов $r_1,\dots,r_{n+1}$ в $\mathbb{R}^{n-1}$ попарно касаются друг друга внешним образом. Доказать, что 
\[
	(n-1)(r_1^{-2} + \dots + r_{n+1}^{-2}) = (r_1^{-1} + \dots + r_{n+1}^{-1})^2.
\]

\item Докажите, что высоты симплекса с вершинами $A_1,\dots,A_{n+1}$ пересекаются в одной точке (такие симплексы называют ортоцентрическими) тогда и только тогда, когда выполнено одно из следующих эквивалентных условий:
  \begin{enumerate}
	  \item При $i \ne j \ne 1$ величина $\langle \overline{A_i A_1}, \overline{A_i A_j} \rangle$ не зависит от $i,j$.
  \item При попарно различных $i,j,k,l$ рёбра $A_i A_j$ и $A_k A_l$ перпендикулярны.
  \item Одна из $(n-1)$-мерных граней -- ортоцентрический симплекс и высота, опущенная на неё, попадает в ортоцентр.
  \item При попарно различных $i,j,k,l$ выполняется соотношение 
  \[
  |A_i A_j|^2 + |A_k A_l|^2 = |A_i A_k|^2 + |A_j A_l|^2 = |A_i A_l|^2 + |A_j A_k|^2.
  \]
  \item Существуют такие числа $a_0,\dots,a_{n+1} \in \mathbb{R}$, что $|A_i A_j|^2 = a_i + a_j$. Доказать, что $a_0^{-1} + \dots + a_{n+1}^{-1} = 0$.
  \item Существует такая точка $A_0$ (ортоцентр), что при попарно различных $i,j,k,l$ выполняются соотношения $(*)$. В этом случае существуют такие числа $a_0,\dots,a_{n+1} \in \mathbb{R}$, что $|A_i A_j|^2 = a_i + a_j$. Доказать, что $a_0^{-1} + \dots + a_{n+1}^{-1} = 0$.
  \end{enumerate}

\item В пространстве размерности $2n+1$ расположен куб (той же размерности) так, что координаты всех его вершин -- целые числа.  
Доказать, что длина ребра куба -- целое число.

\item Шары радиусов $R_1,\dots,R_n$ с центрами $A_1,\dots,A_n$ имеют общую точку. Доказать, что если $B_i B_j \le A_i A_j$, то шары радиусов $R_1,\dots,R_n$ с центрами $B_1,\dots,B_n$ тоже имеют общую точку.
\end{enumerate}

\section{2.2 Полупространства, выпуклые многогранники и полиэдры.}

Выпуклым полиэдром называется множество, образованное пересечением конечного числа полуплоскостей; ограниченные полиэдры называют выпуклыми многогранниками.

\begin{enumerate}
\item Докажите, что $k$-мерное подпространство -- полиэдр.

\item Представить луч $x_i = t$, $1 \le i \le n$, $t \ge 0$ как пересечение полуплоскостей.

\item Найти поляру полуплоскости $\langle a, x \rangle \le b$ в $\mathbb{R}^n$.

\item Найти поляру гиперплоскости $\langle a, x \rangle = b$ в $\mathbb{R}^n$.

\end{enumerate}
Ортогональным дополнением $L^\perp$ подпространства $L$ в $\mathbb{R}^n$ называется совокупность векторов $y \in \mathbb{R}^n$ таких, что $\langle x, y \rangle = 0$ $\forall x \in L$.  
\begin{enumerate}[resume*]
\item Доказать, что если $L$ -- подпространство, то $L^\perp$ совпадает с полярой $L$.
\item Пусть $L_1$ и $L_2$ -- два подпространства в $\mathbb{R}^n$. Как выражается ортогональное дополнение пересечения $L_1$ и $L_2$ через $L_1^\perp$ и $L_2^\perp$?

\item Доказать, что выпуклая оболочка конечного числа точек -- выпуклый многогранник (т.е. пересечение конечного числа полупространств).
\end{enumerate}
Коническая оболочка конечного числа  векторов называется конечнопорожденным конусом.
\begin{enumerate}[resume*]
\item Доказать, что конечнопорожденный конус в $\mathbb{R}^n$ замкнут.

\item Доказать, что конечнопорожденный конус является полиэдром.

\item Доказать, что полярой многогранника, содержащего начало координат внутри себя, является многогранник.

\item Установить взаимно-однозначное соответствие между $k$-мерными гранями многогранника $M$ (содержащего начало координат внутри себя) и $(n-k-1)$-мерными гранями поляры $M$.

\item Пусть $M$ -- многогранник и точка $a$ лежит внутри него. Тогда объем многогранника $V_n(M)$ равен выражению 
\[
n^{-1} \sum_i d(a, M_i) V_{n-1}(M_i),
\]
где $M_i$ -- $(n-1)$-мерные грани $M$.

\item Пусть $M$ -- плоский выпуклый многоугольник, $M_\rho$ -- множество точек $\mathbb{R}^2$, находящихся на расстоянии $\le \rho$ от $M$. Доказать формулу Штейнера:
\[
S(M_\rho) = S(M) + \rho P(M) + \pi \rho^2,
\]
где $S(A)$ -- площадь фигуры $A$, $P(A)$ -- ее периметр.

\end{enumerate}

\section{2.3 Выпуклая геометрия.}

\begin{enumerate}
\item Описать выпуклые множества на прямой.

\item Описать выпуклые множества на плоскости, целиком содержащие некоторую прямую.

\item Описать все разбиения плоскости на 2 выпуклых множества.

\item Выпуклая кривая -- это граница плоского выпуклого множества. Сколько точек пересечения может быть у двух выпуклых кривых?

\item Доказать следующий признак выпуклости: замкнутое множество $X$ выпукло тогда и только тогда, когда для любой точки, принадлежащей $X$, существует ровно одна ближайшая к ней точка $X$.

\item $A$ -- ограниченное замкнутое выпуклое множество в $\mathbb{R}^3$. $\operatorname{extr} A$ -- множество его крайних точек. Привести пример или доказать, что такое $A$ не существует, если:
  \begin{enumerate}
  \item $\operatorname{extr} A = \varnothing$;
  \item $\operatorname{extr} A = \{0\}$;
  \item $\operatorname{extr} A$ -- не замкнуто.
  \end{enumerate}
\newpage
\item Описать следующие множества как решения системы линейных неравенств:
  \begin{enumerate}
  \item круг $x_1^2 + x_2^2 \le 1$;
  \item отрезок с концами $(0,0,0)$ и $(1,1,1)$ в $\mathbb{R}^3$;
  \item выпуклую оболочку параболы $x_2 = x_1^2$ и точки $(0,1)$;
  \item выпуклую оболочку кривой моментов $(t,t^2,\dots,t^d)$ в $\mathbb{R}^d$, $d=2,3,4$.
  \end{enumerate}

\item Найти выпуклую оболочку прямых $\{x: x_3=0, x_2=1\}$ и $\{x: x_1=0, x_2=0\}$ в $\mathbb{R}^3$.

\item Найти поляру:
  \begin{enumerate}
	  \item отрезка $\left[(0,0);(1,1)\right]$;
  \item выпуклой оболочки параболы $x_2 = \frac{x_1^2}{2}$;
  \item прямоугольника $-2 \le x_1 \le 1$; $-3 \le x_2 \le 2$;
  \item множества $A_p = \{x: |x_1|^p + |x_2|^p \le 1\}$, $p \ge 1$.
  \end{enumerate}

\item Провести опорную прямую ко внутренности параболы $x_2 = \frac{x_1^2}{2}$ через точку $(2,2)$.

\item Пусть $A$ -- выпуклое подмножество в $\mathbb{R}^n$. Конус $K$ называется рецессивным для $A$, если $x + K \subset A$ $\forall x \in A$. Найти рецессивные конусы для внутренности параболы $x_2 = \frac{x_1^2}{2}$, внутренности гиперболы $x_1^2 - x_2^2 = 1$, полости гиперболоида $x_1^2 - x_2^2 - x_3^2 = 1$.

\item Пусть $A$ -- непустое подмножество $\mathbb{R}^n$. Тогда каждая точка выпуклой оболочки $A$ есть выпуклая комбинация не более чем $n+1$ точки из $A$ (теорема Каратеодори).

\item Всякое конечное множество из $\mathbb{R}^n$, состоящее не менее чем из $n+2$ точек, можно разбить на два непустых подмножества, выпуклые оболочки которых пересекаются (теорема Радона).

\item Пусть $\{A_i\}_{i=1}^N$ -- семейство замкнутых выпуклых подмножеств $\mathbb{R}^n$, каждые $n+1$ из которых имеют непустое пересечение. Тогда все эти множества имеют общую точку (теорема Хелли).

\item Пусть $A$ -- выпуклое замкнутое множество в $\mathbb{R}^2$. Доказать, что его с любой точностью можно аппроксимировать многогранником $M$.

\item Доказать результат задачи 14 в $\mathbb{R}^n$.

\item Найти крайние точки куба, симплекса, правильного октаэдра.

\item Доказать, что замыкание выпуклого многогранника и его внутренность выпуклы.

\item Пусть $A$ -- непустое выпуклое множество в $\mathbb{R}^n$ и $L$ -- аффинное множество в $\mathbb{R}^n$, не пересекающееся с $A$. Доказать, что существует гиперплоскость, содержащая $L$ и не пересекающаяся с $A$.

\item Пусть $A$ и $B$ -- два замкнутых выпуклых непересекающихся множества, причем одно из них ограничено. Доказать, что существует гиперплоскость, которая их строго разделяет.

\item Пусть $A$ -- замкнутое множество с непустой внутренностью. Если в каждой точке его границы существует опорная гиперплоскость, то $A$ -- выпукло.

\item Пусть $A$ -- замкнутое выпуклое подмножество $\mathbb{R}^n$. Тогда оно совпадает с пересечением всех содержащих его (замкнутых) полупространств.

\item Пусть $A$ -- выпуклое множество, которое можно покрыть полосой ширины единица и не может быть покрыто полосой меньшей ширины. Тогда в $A$ можно вписать круг радиуса $\frac{1}{3}$ (теорема Бланшке).

\end{enumerate}

\section{2.4 Теория линейных неравенств и геометрия.}

В этом параграфе $\mathbb{R}^n$ то же, что и в § 1.4. Множество $K \subset \mathbb{R}^n$ называется выпуклым конусом, если оно является выпуклым множеством, инвариантным относительно гомотетий ($\alpha K \subset K$ $\forall \alpha > 0$).

Конус $K \subset \mathbb{R}^n$ называется конечно порождённым, если он есть коническая оболочка конечного числа векторов.

\begin{enumerate}
\item Доказать, что конечнопорожденный конус замкнут.

\item Доказать, что для того чтобы множество в $\mathbb{R}^n$ было конечнопорожденным конусом, необходимо и достаточно, чтобы оно было решением конечной системы неравенств.
\end{enumerate}

Множество $\Pi \subset \mathbb{R}^n$ называется полиэдральным, если оно есть пересечение конечного числа полупространств.

\begin{enumerate}[resume*]

\item Доказать, что полиэдральное множество, не содержащее прямых, есть сумма многогранника и конечно порождённого конуса.

\item Доказать, что поляра конечно порождённого конуса есть конечно порождённый конус.

\item Если неравенство $\sum\limits_{j=1}^n a_{0j} x_j \le 0$ является следствием системы $\sum\limits_{j=1}^n a_{ij} x_j \le 0$, $1 \le i \le m$, то существуют неотрицательные целые числа $\lambda_1,\dots,\lambda_m$ такие, что $a_{0j} = \sum\limits_{i=1}^m \lambda_i a_{ij}$ (теорема Минковского–Фаркаша).

\item Система $\sum\limits_{j=1}^n a_{ij} x_j - b_i \le 0$, $1 \le i \le m$ несовместна тогда и только тогда, когда существуют такие неотрицательные $\lambda_1,\dots,\lambda_N$, что $\sum\limits_{i=1}^N \lambda_i a_{ij} = 0$, $1 \le j \le n$ и $\sum\limits_{i=1}^m \lambda_i b_i < 0$ (теорема Фань Цзи — А.Д.Александрова).

\item Имеет место альтернатива:  
либо система $a_{i1}x_1 + \dots + a_{in}x_n \le b_i$, $1 \le i \le m$ совместна,  
либо система уравнений  
\begin{gather*}
a_{1j}y_1 + \dots + a_{mj}y_m = 0, 1 \le j \le n \\ 
a_{1}y_1 + \dots + a_{m}y_m = -1  
\end{gather*}
имеет неотрицательное решение.

\item Пусть $K_1$ и $K_2$ — конечно порождённые конусы. Доказать, что  
	$(K_1 + K_2)^{\circ} = K_1^\circ \cap K_2^\circ$ и $(K_1 \cap K_2)^\circ = K_1^\circ + K_2^\circ$.

\end{enumerate}

\section{2.5 Указания, решения и ответы.}

\subsection{2.1 Введение в многомерный мир.}

1. \textit{Ответ:} $16,32,24,8$.

2. Проще всего, пожалуй, задать правильный четырёхмерный симплекс $\Sigma^4$ в $\mathbb{R}^5$:  
$\Sigma^4 = \{ x = (x_1,x_2,x_3,x_4,x_5) \mid x_i \ge 0,\ 1 \le i \le 5,\ \sum x_i = 1 \}$.

3. \textit{Ответ:} $5,10,10,5$.

4. \textit{Ответ:} $8,24,32,16$.

5. \textit{Ответ:} $v_0 - v_1 + v_2 - v_3 = 0$.

6. \textit{Ответ:} рассмотрим симплекс, заданный в задаче 2. Тогда подходит плоскость в $\mathbb{R}^5$, задаваемая уравнениями  
$\sum x_i = 1;$ \\$x_1 + x_2 - 2x_3 = 0;\ 2x_3 - x_4 - x_5 = 0$.

7. \textit{Ответ:} можно двигаться по окружности $(x_1 - 1)^2 + x_2^2 = 1$, $x_3 = x_4 = 0$.  
Она пересекает гиперплоскость $x_4 = 0$ ровно в двух указанных точках.

8. \textit{Ответ:} сделаем параллельные переносы на векторы $(0,0,0,1)$, $(2,0,0,0)$ и $(0,0,0,-1)$, после чего повернём $O_2$ в гиперплоскости $x_4 = 0$ в требуемое положение.

9. \textit{Ответ:} можно. Вначале развернём узел в $\mathbb{R}^3$, отмечая точки, в которых происходят самопересечения. Чтобы получить «правильное» развязывание (без самопересечений), будем приподнимать окрестности одной из пересекающихся нитей вдоль четвёртой координаты, затем опускать обратно в $\mathbb{R}^3$.

10. \textit{Ответ:} $P_1 = \{x : x_1 = x_2 = x_3 = 0\}$; $P_2 = \{x : x_3 = x_4 = 0\}$.

11. \textit{Ответ:} $P_1 = \{x : x_1 = x_2 = 0\}$; $P_2 = \{x : x_2 = x_3 = 1\}$.

12. 
\begin{enumerate}[label=\asbuk*)]
\item Пусть куб задан неравенствами $0 \le x_i \le 1$, а секущая гиперплоскость $x_1 + x_2 + x_3 + x_4 = a$. При $0 < |a-2| < 1$ добавление ещё четырёх плоскостей приводит к тому, что от тетраэдра отрезаются вершины. Полученный многогранник имеет 12 вершин (12 рёбер гиперкуба, пересекаемых данной гиперплоскостью), 18 вершин, 4 треугольника и 4 шестиугольные грани.  

При $a = 2$ такой многогранник вырождается в правильный октаэдр (выпуклая оболочка 6 вершин гиперкуба с $\sum x_i = 2$).

\item Ребро тетраэдра $\{x : \sum x_i = a, x_i \ge 0\}$ равно $\sqrt{2}(a-1)$. Поскольку объём правильного тетраэдра равен $\frac{\sqrt{2}}{12}s^3$, $s$ — длина ребра, то получаем:  
\[
V = \frac{a^3}{3},\ \text{при}\ 0 < a \le 1;\
V = \frac{a^3}{3} - 4\frac{(a-1)^3}{3},\ \text{при}\ 1 < a \le 2.\]  
При $2 < a \le 4$\quad$V(a) = V(4-a)$ из симметрии.
\end{enumerate}
13. При любом $n$. Чтобы показать это, рассмотрим кривую моментов в $\mathbb{R}^4$: $\{(t,t^2,t^3,t^4), t \in \mathbb{R}\}$ и рассмотрим на ней $n$ произвольных точек. Пусть им соответствуют значения параметра $t_1 < t_2 < \dots < t_n$. Покажем, что выпуклая оболочка этих точек даёт многогранник с требуемым свойством. Две вершины связаны ребром, если через них можно провести гиперплоскость так, чтобы все остальные вершины лежали по одну сторону от неё. На кривой моментов линейная функция в $\mathbb{R}^4$ равна многочлену от параметра $t$ степени 4. Многочлен $(t-t_i)(t-t_j)^2$ определяет требуемую гиперплоскость.

16. \textit{Ответ:} $\{ x \in \mathbb{R}^n \mid \sum\limits_{i=1}^n | \frac{x_i}{a_i} | \le 1 \}$.

17. \textit{Ответ:} $\left( \sum\limits_{i=1}^n a_i^2 \right)^{\frac{1}{2}}$.

18. \textit{Ответ:} $\min\limits_i a_i$.

19. \textit{Ответ:} (при $a = (1,1,1)$) $\sqrt{\frac{3}{2}}$.

20. \textit{Ответ:} Если $a_1 \ge a_2 \ge a_3$, то $\sqrt{a_1^2 + a_2^2}$.

28. \textit{Указание:} $V = \int_0^h \left( \frac{t}{h} \right)^{n-1} V_{n-1}\, dt = \frac{1}{n} h^n V_{n-1}$.

29. \textit{Ответ:} $V_n = \frac{\sqrt{n+1}}{n!}$.

\textit{Указание}. Объём $V$ $(n+1)$-мерного симплекса $\sum x_i \le 1$, $x_i \ge 0$ равен $\frac{1}{(n+1)!}$; высота этого симплекса, опущенная на рассматриваемый симплекс, равна $\frac{1}{\sqrt{n+1}}$. С другой стороны, $V = \frac{1}{n} h V_n$.

30. \textit{Ответ:} $\frac{1}{n!} \sqrt{\frac{n+1}{2^n}}$.

\textit{Указание}. В предыдущей задаче найдем объем правильного симплекса с ребром $\sqrt{2}$

34. \textit{Указание}. Центры сфер образуют $n$-мерный симплекс в $\mathbb{R}^{n-1}$, поэтому

$ \det
\begin{pmatrix}
0 & d_{1,2}^2 & \dots & d_{1,n}^2 & 1 \\
d_{2,1}^2 & 0 & \dots & d_{2,n}^2 & 1 \\
\vdots & \vdots & \ddots & \vdots & \vdots \\
d_{n,1}^2 & d_{n,2}^2 & \dots & 0 & 1 \\
1 & 1 & \dots & 1 & 0
\end{pmatrix}
= 0,
$
где $d_{ij} = r_i + r_j$. Разделим $i$-ю строку и $i$-й столбец (кроме последних) на $r_i^2$. Вычтем затем последнюю строку и последний столбец из остальных строк и столбцов. После очевидных сокращений придём к равенству

\[
\det \begin{pmatrix}
	-1 & 1 & 1 & \dots & 1 & r_1^{-1} \\
1 & -1 & 1 & \dots & 1 & r_2^{-1} \\
\vdots & \vdots & \vdots & \ddots & \vdots & \vdots \\
1 & 1 & 1 & \dots & -1 & r_n^{-1} \\
r_1^{-1} & \dots & \dots & \dots & r_{n-1}^{-1} & 0
\end{pmatrix} = 0.
\]

$P\sum r_i^{-2} + 2Q \sum_{i > j} r_i^{-1} r_j^{-1}$, где $P = -\det
\begin{pmatrix}
	-1 & 1 & \dots & 1 \\
	1 & -1 & \dots & 1 \\ 
	\vdots & \vdots & \ddots & \vdots \\
	1 & 1 & \dots & -1
\end{pmatrix},$

$Q = \det
\begin{pmatrix}
	1 & 1 & \dots & 1 \\
	1 & -1 & \dots & 1 \\ 
	\vdots & \vdots & \ddots & \vdots \\
	1 & 1 & \dots & -1
\end{pmatrix}$. (Определители порядка $n$). Прибавив к первой строке $P$ все остальные строки, получим $P = -(n-2)Q$.

35. \textit{Указание}.  
\begin{enumerate}[label=\asbuk*)]
\item Если $A_0$ — ортоцентр, $e_i = \overline{A_1 A_i}$ и $x = A_1 A_0$, то $\langle x, e_i - e_j \rangle = 0$ и $\langle x - e_i, e_j \rangle = 0$. Следовательно, $\langle x, e_i \rangle = \text{const}$ и $\langle e_i, e_j \rangle = \langle x, e_i \rangle$.  

Если $\langle e_i, e_j \rangle = \text{const}$, то существуют такие векторы $\varepsilon_i,$ что $\langle \varepsilon_i, \varepsilon_j \rangle = \delta_{ij}.$ Положим $x = c(\varepsilon_2 + \dots + \varepsilon_{n+1})$, где $c = \langle e_i, e_j \rangle$. Тогда $ \langle x, e_i - e_j \rangle = 0$ и $\langle x - e_i, e_j \rangle = 0$. 

\item Равенства $\langle e_i, e_j - e_k \rangle = 0$ и $\langle e_i - e_j, e_k - e_l \rangle = 0$ эквивалентны тому, что $\langle e_i, e_j \rangle = \text{const}$. 

\item Пусть $A_1 H$ — высота, опущенная на грань $A_2 \dots A_{n+1}$. Тогда $A_1 H \perp A_i A_j$ при $i,j \ge 2$. Поэтому $A_k H \perp A_i A_j \Leftrightarrow A_1 A_k \perp A_i A_j$. 

\item Пусть $a,b,c,d$ — векторы сторон четырёхугольника $A_1 A_2 A_3 A_4$. Равенство $a^2 + c^2 = b^2 + d^2$ эквивалентно тому, что $\langle a,c \rangle = \langle b, d \rangle = - \langle b, a + b + c \rangle$, т.е. $\langle a + b, c + d \rangle = 0$.  

\item Если $|A_i A_j|^2 = a_i + a_j$, то $|A_i A_j|^2 + |A_k A_l|^2 = a_i + a_j + a_k + a_l = |A_i A_k|^2 + |A_j A_l|^2$. А если симплекс ортоцентрический, то положим $a_l = \frac{1}{2} (|A_l A_j|^2 + |A_l A_k|^2 - |A_j A_k|^2)$ (эта величина не зависит от $i,j$ и $k$). 

\item Если $i,j,k$ попарно различны и отличны от нуля, то $A_0 A_i \perp A_j A_k$. Равенство $a_0^{-1} + \dots + a_{n+1}^{-1} = 0$ докажем по индукции. Пусть $A_0'$ и $A_0$ — ортоцентры симплексов $A_1 \dots A_n$ и $A_1 \dots A_{n+1}$, $A_0' A_0 = x$ и $A_0 A_{n+1} = y$. Тогда $|A_i A_0|^2 = a_i + a_0' + x^2$ и $|A_i A_{n+1}|^2 = a_i + a_0' + (x+y)^2$, т.е. $a_0 = a_0' + x^2$ и $a_{n+1} = a_0' + (x+y)^2$. Кроме того, $y^2 = a_0 + a_{n+1}$, т.е. $a_0' = -x (x+y)$. Следовательно $a_0^{-1} + \dots + a_{n+1}^{-1} = \frac{1}{-x y} + \frac{1}{y(x+y)} - \frac{1}{x(x+y)} = 0$. База индукции: точка $A_0$, лежащая на отрезке $A_1 A_2$.
\end{enumerate}

36. \textit{Указание}. Объём рассматриваемого куба и квадрат длины ребра — целые числа. Следовательно, $a = \frac{a^{2n+1}}{(a^2)^n} $ — целое число.

37. \textit{Указание}. Пусть $P$ — общая точка исходных шаров, $f(x) = \max\limits_i \frac{X B_i}{P A_i}$, $f(X) \to \infty$ при $X \to \infty$, поэтому $f$ имеет точку минимума $Q$. Достаточно доказать, что $f(Q) \le 1$ (тогда $Q B_i \le P A_i \le R_i$). График функции $y = f(X)$ — граница пересечения конусов с вершинами $B_1,\dots,B_n$. Точка минимума принадлежит границам не менее чем двух конусов с вершинами $B_i$ и $B_j$. Рассмотрев ограничение $f(X)$ на прямую $B_i B_j$, легко проверить, что $f(Q) \le 1$.

\subsection{2.3 Выпуклая геометрия.}

1. \textit{Ответ:} интервалы, отрезки, открытые и замкнутые лучи, точка и вся прямая.

2. \textit{Ответ:} внутренность некоторой полосы $a < x_2 < b$ и выпуклые множества на прямых $x_2 = a$ и $x_2 = b$.

3. Одним из множеств является замкнутая полуплоскость или полуплоскость с лучом на границе.  

\textit{Указание}. Если выпуклое множество содержится в угле, меньшем $180^\circ$, то дополнение к нему — невыпукло. Значит, все отделяющие прямые — параллельны.

4. \textit{Ответ:} $0,1,\dots,\infty$.

5. \textit{Указание.} Обозначим через $x(y)$ ближайшую к $y$ точку из $X$. Достаточно проверить, что гиперплоскость $\langle y - x(y), x \rangle = \langle y - x(y), x(y) \rangle$ строго отделяет $y$ от $X$, и что $X$ есть пересечение полупространств $\langle y - x(y), x \rangle \le \langle y - x(y), x(y) \rangle$.

6. \textit{Ответы:}  
\begin{enumerate}[label=\asbuk*)]
	\item не существует (теорема Крейна–Мильмана);  
	\item $\{0\}$;  
	\item $\operatorname{Conv}\left(\{x : |x|^2 \le 1\}, (-1,0,1), (1,0,1)\right)$.
\end{enumerate}

7. \textit{Ответы:} 
\begin{enumerate}[label=\asbuk*)]
	\item $\{x : \langle x,y \rangle \le 1, y_1^2 + y_2^2 = 1\}$;  
	\item $\{x : x_1 \le x_2 \le x_3 \le x_4;\ 0 \le x_i \le 1\}$;  
	\item $\{x : x_2 - 2tx_1 + r^2 \ge 0,\ |t| \ge 1\}$;  
	\item
$d=2$ — $\{x : x_2 - 2tx_1 + t^2 \ge 0\}$,

$d=3$ — $\{x : x_2 - 2tx_1 + t^2 \ge 0\}$,  

$d=4$ — $\{x : x_4 + 2(pt)x_3 + (t^2 - 4pt + q)x_2 + 2t(pt - q)x_1 + t^2 q^2 \ge 0\}$.  
\end{enumerate}
\textit{Указание}. Линейная функция в $\mathbb{R}^4$ на кривой моментов является многочленом параметра $t$, а многочлен 3-й степени имеет хотя бы один вещественный корень. Найти многочлены 4-й степени, принимающие неотрицательные значения и имеющие вещественный корень.

8. \textit{Ответ:} объединение этих прямых и полосы $0 < x_2 < 1$. 

\textit{Указание}. Взяв произвольную точку полосы, найдём пересечение плоскости, содержащей одну из прямых и выбранную точку, с другой прямой. Осталось убедиться, что точка пересечения расположена дальше от выбранной прямой, чем выбранная точка полосы.

\newpage
9. \textit{Ответы:}  
\begin{enumerate}[label=\asbuk*)]
\item $\{y : y_1 + y_2 \le 1\}$;  
\item $\{y : -y_2 \ge \frac{y_1^2}{2\}}$;  
\item $\operatorname{Conv}\left((0,1/2),(0,-1/3),(1,0),(-1/2,0)\right)$.
\end{enumerate}
\textit{Указание к} в). непосредственно проверяется, что указанные точки принадлежат поляре. Записав их выпуклую оболочку в виде неравенств $\langle a, x \rangle \le 1$, заметим, что векторы $a$ в этих неравенствах совпадают с вершинами прямоугольника.
\begin{enumerate}[resume*]
	\item $A_p^\circ = A_{p/{p-1}}$, $A_1^\circ = A_\infty$.
\end{enumerate}
\textit{Указание}. Воспользоваться неравенством Гёльдера 
\[
|x_1 y_1 + x_2 y_2| \le (x_1^p + x_2^p)^{\frac{1}{p}} (y_1^q + y_2^q)^{\frac{1}{q}}, \quad \frac{1}{p} + \frac{1}{q} = 1,
\]
где равенство достигается лишь при $x_1^p = y_1^p$ и $x_2^p = y_2^p$.
\end{document}
